%! Sinusoidal Function Context and Data Modeling — Unit 6, Lesson 6.8 — Group Activity (Answer Key)
\documentclass[10pt]{article}
\usepackage{precalculus-article}
\usepackage{precalculus-key}   % pulls in -boxes; do NOT also load -boxes
\newcommand{\TallMath}[1]{$\displaystyle #1\rule[-1.4em]{0pt}{3.2em}$}

\begin{document}

\pageheader{Unit 6, Lesson 6.8}{Group Activity --- Answer Key}
\namepartnerperiod

\vspace{0.10in}

%----------------------------------------------------------------------
% Context
%----------------------------------------------------------------------
\begin{scenariobox}[Context: Modeling Periodic Real-World Data]{plumlight}
Sinusoidal functions appear in temperature records, tidal data, daylight hours,
and many other natural phenomena. In this activity your group will build and
evaluate a sinusoidal model from real-world data. Choose \textbf{one tier}
based on your readiness. Everyone should be ready to explain the reasoning
behind each parameter your group chooses.
\end{scenariobox}

\vspace{0.08in}

%----------------------------------------------------------------------
% Tier R Key
%----------------------------------------------------------------------
\begin{tcolorbox}[colback=white, colframe=black!40,
    title=\textbf{Tier R --- Reinforce (Key)},
    fonttitle=\bfseries, arc=2mm, left=3mm, right=3mm, top=2mm, bottom=2mm]

\textbf{Context:} The table gives average monthly temperature (${}^\circ$F).
Max = 80 (month 7), Min = 46 (month 12/Jan area).

\vspace{4pt}
\begin{center}
\renewcommand{\arraystretch}{1.35}
\small
\begin{tabular}{|c|c|c|c|c|c|c|c|c|c|c|c|c|}
\hline
Month $x$ & 1 & 2 & 3 & 4 & 5 & 6 & 7 & 8 & 9 & 10 & 11 & 12 \\
\hline
Temp $T$ & 48 & 51 & 58 & 65 & 72 & 78 & \textbf{80} & 79 & 72 & 63 & 54 & \textbf{46} \\
\hline
\end{tabular}
\end{center}

\vspace{4pt}
\textbf{Step 1.} $A = \dfrac{80-46}{2} = $ \ans{17} \qquad $D = \dfrac{80+46}{2} = $ \ans{63}

\vspace{4pt}
\textbf{Step 2.} Next max at month \ans{19}, so $k = $ \ans{12}. $B = \dfrac{2\pi}{12} = $ \ans{$\tfrac{\pi}{6} \approx 0.524$}

\vspace{4pt}
\textbf{Step 3.} $C = x_{\max} = $ \ans{7}

\vspace{4pt}
\textbf{Step 4.} $T(x) = $ \ans{$17\cos\!\left(\tfrac{\pi}{6}(x-7)\right)+63$}

\vspace{4pt}
\textbf{Step 5.} $T(1) = 17\cos\!\left(\tfrac{\pi}{6}(1-7)\right)+63 = 17\cos(-\pi)+63 = 17(-1)+63 = $ \ans{46}. Residual: \ans{$48-46=2$; fairly close.}

\vspace{4pt}
\textbf{Step 6.} $T(4.5) \approx 17\cos\!\left(\tfrac{\pi}{6}(4.5-7)\right)+63 = 17\cos\!\left(-\tfrac{5\pi}{12}\right)+63 \approx 17(0.259)+63 \approx $ \ans{$67.4^\circ$F}

\end{tcolorbox}

\vspace{0.08in}

%----------------------------------------------------------------------
% Tier A Key
%----------------------------------------------------------------------
\begin{tcolorbox}[colback=white, colframe=black!40,
    title=\textbf{Tier A --- Apply (Key)},
    fonttitle=\bfseries, arc=2mm, left=3mm, right=3mm, top=2mm, bottom=2mm]

\textbf{Context:} Tidal height (ft) over 12 hours.

\vspace{4pt}
\begin{center}
\renewcommand{\arraystretch}{1.3}
\small
\begin{tabular}{|c|c|c|c|c|c|c|c|c|c|c|c|c|c|}
\hline
$t$ & 0 & 1 & 2 & 3 & 4 & 5 & 6 & 7 & 8 & 9 & 10 & 11 & 12 \\
\hline
$h$ & 1.5 & 3.8 & 5.5 & 6.2 & 5.5 & 3.8 & 1.5 & $-0.8$ & $-2.5$ & $-3.2$ & $-2.5$ & $-0.8$ & 1.5 \\
\hline
\end{tabular}
\end{center}

\vspace{4pt}
\textbf{1.} Max $= $ \ans{6.2} at $t = $ \ans{3} hr. \quad Min $= $ \ans{$-3.2$} at $t = $ \ans{9} hr.

\vspace{4pt}
\textbf{2.} $A = \dfrac{6.2-(-3.2)}{2} = $ \ans{4.7} \quad $D = \dfrac{6.2+(-3.2)}{2} = $ \ans{1.5}

Midline context: \ans{The average tide height above/below mean sea level --- approximately where the water would be with no tidal variation.}

\vspace{4pt}
\textbf{3.} The pattern repeats every 12 hours (max at $t=3$, next max at $t=15$).

$k = $ \ans{12} hr \quad $B = \dfrac{2\pi}{12} = $ \ans{$\tfrac{\pi}{6} \approx 0.524$}

\vspace{4pt}
\textbf{4.} $h(t) = $ \ans{$4.7\cos\!\left(\tfrac{\pi}{6}(t-3)\right)+1.5$}

\vspace{4pt}
\textbf{5.} $h(0) = 4.7\cos\!\left(\tfrac{\pi}{6}(0-3)\right)+1.5 = 4.7\cos\!\left(-\tfrac{\pi}{2}\right)+1.5 = 4.7(0)+1.5 = $ \ans{1.5 ft}. Residual: \ans{0 (exact match).}

\vspace{4pt}
\textbf{6.} Need $h(t) \geq 2$. Setting $4.7\cos\!\left(\tfrac{\pi}{6}(t-3)\right)+1.5 = 2$ gives $\cos = 0.106$, so
$\tfrac{\pi}{6}(t-3) \approx \pm 1.464$, giving $t \approx 3 \pm 2.8$.
\ans{Entry is safe approximately on $0 \leq t \lesssim 5.8$ and $0.2 \leq t \lesssim 5.8$ --- roughly from $t=0.2$ hr to $t=5.8$ hr, and again from $t=9$ to $t=12$ hr (when the tide exceeds 2 ft again near the end of the cycle). Students may also express this as approximately hours 0--6 and 9--12.}

\begin{teachernote}
Exact computation requires inverse cosine. Accept reasonable approximations read from the data table: tide exceeds 2 ft from about $t = 0$ to $t = 6$ (reading from the table). Tier E students may go further.
\end{teachernote}

\end{tcolorbox}

\vspace{0.08in}

%----------------------------------------------------------------------
% Tier E Key
%----------------------------------------------------------------------
\begin{tcolorbox}[colback=white, colframe=black!40,
    title=\textbf{Tier E --- Extend (Key)},
    fonttitle=\bfseries, arc=2mm, left=3mm, right=3mm, top=2mm, bottom=2mm]

\textbf{1.} City A: max = 75 (mo. 8), min = 58 (mo. 1).
$A_A = \tfrac{75-58}{2} = $ \ans{8.5},
$D_A = \tfrac{75+58}{2} = $ \ans{66.5},
$k_A = 12$, $B_A = \tfrac{\pi}{6}$, $C_A = 8$.
Model: \ans{$T_A(x) = 8.5\cos\!\left(\tfrac{\pi}{6}(x-8)\right)+66.5$}

City B: max = 94 (mo. 7), min = 28 (mo. 1).
$A_B = \tfrac{94-28}{2} = $ \ans{33},
$D_B = \tfrac{94+28}{2} = $ \ans{61},
$k_B = 12$, $B_B = \tfrac{\pi}{6}$, $C_B = 7$.
Model: \ans{$T_B(x) = 33\cos\!\left(\tfrac{\pi}{6}(x-7)\right)+61$}

\vspace{4pt}
\textbf{2.} City B has the larger amplitude (\ans{33 vs.\ 8.5}).
\ans{City B experiences far greater seasonal temperature swings --- it gets much hotter in summer and much colder in winter than City A, which has a mild coastal climate.}

\vspace{4pt}
\textbf{3.} City A: $T_A(1) = 8.5\cos\!\left(-\tfrac{7\pi}{6}\right)+66.5 \approx 8.5(-0.866)+66.5 \approx $ \ans{59.1}. Residual: $58 - 59.1 = $ \ans{$-1.1$}.

City B: $T_B(1) = 33\cos\!\left(-\pi\right)+61 = 33(-1)+61 = $ \ans{28}. Residual: $28 - 28 = $ \ans{0 (exact)}.

\vspace{4pt}
\textbf{4.} City B residuals: $T_B(4) = 33\cos\!\left(-\tfrac{\pi}{2}\right)+61 = 61$; actual 66; residual \ans{$+5$}. $T_B(7) = 33(1)+61 = 94$; actual 94; residual \ans{0}. $T_B(10) = 33\cos\!\left(\pi\right)+61 = 28$; actual 60; residual \ans{$+32$}. Large residual at $x=10$ suggests the sinusoidal model is less accurate for fall months; \ans{the actual fall temperatures stay warmer than the model predicts, perhaps due to thermal lag.}

\vspace{4pt}
\textbf{5.} Contextual domain: \ans{$1 \leq x \leq 12$ (one year) for each city.} City B at $x=14$: $T_B(14) = 33\cos\!\left(\tfrac{\pi}{6}(14-7)\right)+61 = 33\cos\!\left(\tfrac{7\pi}{6}\right)+61 \approx 33(-0.866)+61 \approx $ \ans{32.4$^\circ$F}.
\ans{This is outside the 12-month contextual domain, so the prediction may be less trustworthy --- especially since real temperatures vary year to year and the large residuals in fall suggest the model doesn't perfectly capture City B's seasonal pattern.}

\end{tcolorbox}

\end{document}
